%% ma: L12-B14-0025
%% lop: 12
%% chuong: 5
%% bai: 14
%% dang: 12.14.16
%% loai: TN4
%% muc: TH
%% dap_an: "B"
%% nguon: "CĐ2.5.pdf giáo viên gửi 27/10/2026 (M07, Lê Minh Khôi – Nguyễn Anh Khôi, Chương II Chủ đề 5 Tích có hướng của hai vectơ), Câu 9"
%% kien_thuc: "Bài 14 (tích có hướng, phương trình mặt phẳng qua ba điểm)"
%% trang_thai: cho-duyet
%% ly_do: "Dạng mới 12.14.16 đề xuất, chờ giáo viên duyệt"
\begin{ex}
Trong không gian $Oxyz$, cho ba điểm $A(1;1;1)$, $B(-1;7;-3)$ và $C(2;1;0)$. Tọa độ của điểm $D$ thuộc trục $Oz$ sao cho bốn điểm $A$, $B$, $C$, $D$ đồng phẳng là
\choice{$D(0;0;-2)$}{\True $D(0;0;3)$}{$D(0;0;-3)$}{$D(0;0;2)$}
\loigiai{$\left[\overrightarrow{AB},\overrightarrow{AC}\right]=(-6;-6;-6)$ nên $(ABC)$: $x+y+z-3=0$.

$D(0;0;z)\in(ABC)\Leftrightarrow z=3$.}
\end{ex}
